\subsection{Complete convex cones in weak spaces}

Lemma 1. --- Let E be a Hausdorff weak space and C a proper cone with vertex 0 in E, that is complete for the uniform structure induced by that of E. Every continuous linear form in E is then the difference between two continuous linear forms in E that are positive in C.

Let \(E'\) be the dual of E and F be the algebraic dual of \(E'\) , with the topology \(\sigma(F, E')\) . Let \(H = C^{\circ} - C^{\circ}\) be the vector subspace of \(E'\) formed by the differences of linear forms that are continuous in E and positive in C (II, p.~44, prop. 4). It is sufficient to show that the orthogonal to H in F is \(\{0\}\) (II, p.~41, Example 1). Then let \(a \in F\) be orthogonal to H; as a is orthogonal to \(C^{\circ}\) , it must belong to the bipolar of C in F. But E is identifiable as a subspace of F, and since C is complete, thus closed in F, we have \(a \in C\) (II, p.~44, th. 1). Similarly a is orthogonal to \(-C^{\circ}\) and therefore \(a \in -C\) . As C is proper, we have a = 0.

PROPOSITION 11. --- Let E be a Hausdorff weak space, and C be a proper convex cone with vertex 0 in E and which is complete in the uniform structure induced by that of E.

Then there exists a set I and a continuous linear mapping u of E in the product space \(R^{I}\) with the following properties :

\begin{enumerate}
\def\labelenumi{\alph{enumi})}
\item
  \(u\) is an isomorphism of \(\mathbf{C}\) on \(u(\mathbf{C})\) for the uniform structures induced respectively by those of \(\mathbf{E}\) and of \(\mathbf{R}^{\mathrm{I}}\) .
\item
  We have \(u(\mathbf{C}) \subset \mathbf{R}_{+}^{\mathrm{I}}\) .
\end{enumerate}

Further, if the uniform structure induced on C by that of E is metrisable, then we can take I = N.

Let \((f_{\iota})_{\iota \in I}\) be a family of continuous linear forms in E such that the finite sums of pseudometrics of the form \((x, y) \mapsto |f_{\iota}(x - y)|\) on \(\mathbf{C} \times \mathbf{C}\) define the uniform structure of C. (If the structure is metrisable we can take \(\mathbf{I} = \mathbf{N}\) .) By lemma 1 we can suppose further that each of the \(f_{\iota}\) is positive in C. Let \(u\) be the linear mapping \(x \mapsto (f_{\iota}(x))_{\iota \in I}\) of E in \(\mathbf{R}^{\mathrm{I}}\) . It is clear that \(u\) is continuous and that \(u(\mathbf{C}) \subset \mathbf{R}_{+}^{\mathrm{I}}\) . The restriction \(u|\mathbf{C}\) is a uniformly continuous mapping that is surjective from C on \(u(\mathbf{C})\) . Further if \(x, y\) in C are such that \(f_{\iota}(x) = f_{\iota}(y)\) for all \(\iota \in I\) , then \(x = y\) since the uniform structure of C is Hausdorff; therefore \(u|\mathbf{C}\) is bijective. Finally, if W is an entourage of the uniform structure of C, then there exists a finite set J of I and a number \(\varepsilon > 0\) such that the relations \(|f_{\iota}(x) - f_{\iota}(y)| \leqslant \varepsilon\) for \(\iota \in J\) imply \((x, y) \in W\) ; therefore \(u|\mathbf{C}\) is an isomorphism of C on \(u(\mathbf{C})\) for the uniform structures being considered.

COROLLARY 1. --- Let E be a Hausdorff weak space and C a proper convex cone of vertex 0 in E that is complete for the uniform structure induced by that of E. Then the mapping \((x, y) \mapsto x + y\) of \(C \times C\) in C is proper.

Because of prop. 11, we can suppose that \(\mathbf{E} = \mathbf{R}^{\mathbf{I}}\) and that \(\mathbf{C} = \mathbf{R}_{+}^{\mathbf{I}}\) (GT, I, § 10.1, cor. 1 and 4). But then the mapping \((x, y) \mapsto x + y\) of \(\mathbf{C} \times \mathbf{C}\) in \(\mathbf{C}\) is written as \(((\xi_{\mathfrak{i}}), (\eta_{\mathfrak{i}})) \mapsto (\xi_{\mathfrak{i}} + \eta_{\mathfrak{i}})\) , and we can restrict ourselves to proving that the continuous mapping \(f: (\xi, \eta) \mapsto \xi + \eta\) of \(\mathbf{R}_{+} \times \mathbf{R}_{+}\) in \(\mathbf{R}_{+}\) , is proper (GT, I, § 10.1, cor.3), Now, for all \(\zeta \in \mathbf{R}_{+}\) , we see that \(f(\zeta)\) is the set of pairs \((\xi, \zeta - \xi)\) such that \(0 \leqslant \xi \leqslant \zeta\) , therefore the inverse image by \(f\) of the interval \([0, \zeta]\) is the set of the \((\xi, \eta) \in \mathbf{R}_{+} \times \mathbf{R}_{+}\) such that \(\xi + \eta \leqslant \zeta\) , which is compact. The conclusion follows applying (GT, I, § 10.3, prop. 7).

COROLLARY 2. --- Let E be a Hausdorff weak space, and C a proper convex cone with vertex 0 in E, that is complete for the uniform structure induced by that of E.

\begin{enumerate}
\def\labelenumi{(\roman{enumi})}
\item
  For every point \(a\) of \(\mathbf{E}\) , the intersection \(\mathbf{C} \cap (a - \mathbf{C})\) is compact.
\item
  Let \(\mathbf{A},\mathbf{B}\) be two closed sets in C. Then \(\mathbf{A} + \mathbf{B}\) is a closed set in C.
\item
  The set of the \((x, y) \in \mathbf{C} \times \mathbf{C}\) such that \(x + y = a\) is compact from cor. 1 and from GT, I, § 10.2, th. 1, b). Now this set is also the set of the \((x, a - x)\) for \(x \in \mathbf{C} \cap (a - \mathbf{C})\) , which proves (i).
\item
  If A and B are closed in C, then \(A \times B\) is closed in \(C \times C\) , therefore \(A + B\) is closed in C after cor. 1 and GT, I, § 10.1, prop. 1.
\end{enumerate}

